\section{Lecture 20: Normed Analytic Rings}
\label{lec:20}

Having defined analytic stacks in \cref{lec:19}, Clausen spends the first part
of the lecture revisiting two points from Scholze's exposition -- the functor
from light condensed anima to analytic stacks, and the description of maps into
(the analytic stack attached to) a compact Hausdorff space -- filling in the
argument that a finite-dimensional compact Hausdorff space is recovered
Tannakially from its category of sheaves. The rest of the lecture develops the
notion of a \emph{norm} on an analytic ring (\cref{def:19:norm}). Clausen gives
an axiomatisation that differs slightly from Scholze's: one of Scholze's
multiplicativity conditions is rephrased, and a new condition (Axiom~4 below)
is added, placing the ``basic disk'' $\mathrm{Spec}(P)$ between the open and the
closed unit disk. Whether Axiom~4 follows from the others is not known in
general (it does over a base that solidifies to $0$). The lecture closes with a
classification: giving a norm together with a topologically-small element $q$ is
the same as a map to $\mathrm{Spec}(\mathbb{Z}[\widehat q^{\pm1}]_\gasx)$ times a
value $N(q)\in(0,1)$; the gaseous base ring of \cref{lec:14} is exactly what
controls norms.

Throughout, ``ring'' means commutative ring and everything is (light)
condensed.

\subsection{The functor from condensed anima}
\label{subsec:20:condensed-anima}

Recall (\cref{ex:19:condensed-anima}) the functor
\begin{equation}
\label{eq:20:condan-to-anstk}
\CondAn^{\mathrm{light}}\ \longrightarrow\ \AnStk .
\end{equation}
It is induced on the generators -- light profinite sets $T\in
\ProFin_{\mathbb N}(\mathrm{Fin})$ -- by
\begin{equation}
\label{eq:20:T-to-spec}
T\ \longmapsto\ \AnSpec\bigl(\Cont(T,\mathbb Z)\bigr),
\end{equation}
where $\Cont(T,\mathbb Z)$ is the \emph{discrete} ring of continuous
($=$ locally constant) $\mathbb Z$-valued functions on $T$, endowed with the
maximal (trivial, ``uncompleted'') analytic ring structure -- i.e.\ all
condensed modules, or all derived condensed modules, over that discrete ring are
declared complete.

Two facts make this functor well-behaved.
\begin{itemize}
\item The assignment \eqref{eq:20:T-to-spec} sends hypercovers of light
profinite sets to $!$-hypercovers satisfying $!$-descent (equivalently
$!$-codescent, a purely formal $\PrL$-fact of Lurie, independent of the shape
of the diagram; \cref{thm:18:tfae,cor:17:descent-codescent}).
\item It preserves finite limits, in particular pullbacks. A pullback of light
profinite sets is a filtered inverse limit of pullbacks of finite sets, which
goes to a filtered colimit of the corresponding situation for finite sets;
the finite case is clear.
\end{itemize}
Consequently the induced functor \eqref{eq:20:condan-to-anstk} is
colimit-preserving and finite-limit-preserving.

\begin{remark}[Algebro-geometric analogue]
\label{rmk:20:gregoric}
Gregori\'c \cite{gregoric2024stonedualitycondensedmathematics} studies the
analogue of \eqref{eq:20:condan-to-anstk} with analytic stacks replaced by fpqc
sheaves in ordinary algebraic geometry over a field, developing a Stone-type
dictionary between condensed mathematics and algebraic geometry in some detail;
Clausen recommends it.
\end{remark}

\begin{remark}[No underlying topological space]
\label{rmk:20:no-underlying-space}
Asked how to recover a topological space from an analytic stack, Clausen was
non-committal: there is an underlying light condensed set, and one could map
that to topological spaces, but there is no evident well-defined ``underlying
topological space'' functor on arbitrary analytic stacks. One can try to extract
one using monomorphisms (open substacks) of analytic stacks; there are several
inequivalent attempts. For a \emph{normed} analytic ring (\S\ref{subsec:20:norms}
below) there is, however, a good way of extracting underlying topological spaces
for objects over it.
\end{remark}

\subsection{Maps into a compact Hausdorff space}
\label{subsec:20:compact-hausdorff}

Let $K$ be a compact Hausdorff space that is metrizable and finite-dimensional
-- equivalently, $K$ embeds as a closed subset of a finite power of the closed
unit interval,
\begin{equation*}
K\ \hookrightarrow\ [0,1]^N .
\end{equation*}
Viewing $K$ as a light condensed set (hence a light condensed anima), we get an
associated analytic stack, abusively also denoted $K$. Scholze made two claims
about it (\cref{lem:19:tanaka}); we reprove them.

\begin{theorem}[Sheaf description and Tannakian points]
\label{thm:20:compact-hausdorff}
Let $K$ be metrizable, finite-dimensional compact Hausdorff.
\begin{enumerate}
\item The derived category of the associated analytic stack is sheaves on the
topological space $K$ valued in the derived category of the trivial base ring:
\begin{equation*}
D(K)\ =\ \Shv\bigl(K;\ D^{\mathrm{cond}}(\mathbb Z)\bigr),
\qquad D^{\mathrm{cond}}(\mathbb Z)=D(\CondAbl),
\end{equation*}
and after base change to an analytic ring $R$ one replaces $D^{\mathrm{cond}}
(\mathbb Z)$ by $D(R)$.
\item For an analytic ring $R$, maps of analytic stacks $\mathrm{Spec}(R)\to K$
are the same as symmetric monoidal, colimit-preserving, $D(\mathbb Z)$-linear
functors
\begin{equation}
\label{eq:20:tannaka}
F\colon\ \Shv\bigl(K;\ D(\mathbb Z)\bigr)\ \longrightarrow\ D(R)
\end{equation}
such that, $!$-locally on $\mathrm{Spec}(R)$, one has
$F\bigl(\Shv(K;D(\mathbb Z)_{\ge0})\bigr)\subseteq D(R)_{\ge0}$.
\end{enumerate}
Both sides of \eqref{eq:20:tannaka} are sets, not merely anima.
\end{theorem}

Here $\Shv(K;D(\mathbb Z))$ uses \emph{discrete} coefficients (the ordinary
derived category of abelian groups). The finite-dimensionality of $K$ enters
only through part~(2), and only in comparing $D(\mathbb Z)$-valued with
$D^{\mathrm{cond}}(\mathbb Z)$-valued sheaves; part~(1) and the reduction to
idempotent algebras below work for an arbitrary compact Hausdorff $K$.

\begin{remark}[Both sides are sets]
\label{rmk:20:both-sets}
The left-hand side of \eqref{eq:20:tannaka} is a set because the diagonal
$K\xrightarrow{\Delta}K\times K$ is a monomorphism in $\CondAn^{\mathrm{light}}$,
and \eqref{eq:20:condan-to-anstk} preserves this pullback; so two maps
$\mathrm{Spec}(R)\to K$ agree in at most one way. The right-hand side is a set
because of $D(\mathbb Z)$-linearity: such an $F$ is the same as a
colimit-preserving functor out of $\Shv(K;\mathrm{Anima})$, which is generated
under colimits by the representables, hence determined on those.
\end{remark}

\begin{remark}[Reduction to idempotent algebras]
\label{rmk:20:idempotents}
The functor $F$ in \eqref{eq:20:tannaka} is determined by where it sends the
constant sheaves $\mathbb Z_Z=i_{Z!}\mathbb Z$ on closed subsets $Z\subseteq K$;
each such must go to an idempotent commutative algebra in $D(R)$. Thus the
right-hand side is
\begin{equation*}
\left\{
\begin{aligned}
&\text{maps }\{\text{closed subsets of }K\}\ \longrightarrow\ \{\text{idempotent
algebras in }D(R)\}\\
&\text{sending filtered }\textstyle\bigcap\text{ to filtered }\varinjlim,\
\text{and }\bigcup\text{ to the Mayer--Vietoris construction}
\end{aligned}
\right\}
\end{equation*}
together with the connectivity condition. The idempotent algebras in $D(R)$ form
a poset (a priori an $\infty$-category, but one checks it is a poset), so a map
of posets satisfying these conditions is unique if it exists. The identification
of $\mathbb Z_Z$ with the ``derived take-away of $K$'' near $Z$ is where
finite-dimensionality is used.
\end{remark}

We indicate the two directions.

\emph{The forward map.} Given $f\colon\mathrm{Spec}(R)\to K$ one takes $f\mapsto
f^*$. When $K$ is profinite, a map $\mathrm{Spec}(R)\to K$ is by construction the
same as a ring map $\Cont(K,\mathbb Z)\to\ur R$ (a filtered colimit of the
finite case $\Cont(K_n,\mathbb Z)$), and there the connectivity condition is
automatic: the unit $\mathbb Z_K$ goes to the unit of $D(R)$, hence to something
connective; every clopen subset then goes to a retract of the unit, and clopen
subsets generate the connective part under colimits, so everything in
$\Shv(K;D(\mathbb Z)_{\ge0})$ lands in $D(R)_{\ge0}$.

\emph{The converse.} Given $F$, choose a surjection $T_0\twoheadrightarrow K$
from a profinite set and form its \v Cech nerve $T_\bullet$ (each
$T_i\subseteq T_0^{\times(i+1)}$ is again profinite, being closed in a product of
profinite sets). Applying $F$ to the pushforwards $\pi_{i*}\mathbb Z$ gives a
cosimplicial diagram of commutative algebras in $D(R)_{\ge0}$; by a K\"unneth
computation -- and because the maps $\pi_i$ are proper (any map of compact
Hausdorff spaces is proper), so proper base change applies -- the whole diagram
is the \v Cech nerve of a single map, and $R\pi_{0*}\mathbb Z$ is concentrated in
degree $0$. Set
\begin{equation*}
R'\ :=\ \Bigl(F\bigl((\pi_0)_*\mathbb Z\bigr),\ \Mod_{F((\pi_0)_*\mathbb Z)}
\bigl(D(R)\bigr)\Bigr),
\end{equation*}
the analytic ring on the algebra $F((\pi_0)_*\mathbb Z)\in\CAlg(D(R)_{\ge0})$
with the induced structure. One produces a map $\mathrm{Spec}(R')\to
\mathrm{Spec}(R)$ (proper: only the underlying ring changes) together with a map
from its \v Cech nerve to that of $T_0\to K$, and hence, by descent, a map
$\mathrm{Spec}(R)\to K$ -- \emph{provided} $\mathrm{Spec}(R')\to\mathrm{Spec}(R)$
is a $!$-cover.

\begin{remark}[The descendability reduction]
\label{rmk:20:descendability}
As $\mathrm{Spec}(R')\to\mathrm{Spec}(R)$ is proper, \cref{prop:18:check}(2a)
reduces $!$-cover to the finitary membership $\ur R\in\langle\operatorname{im}
(D(R')\to D(R))\rangle$. Applying $F$, it suffices that
\begin{equation*}
\mathbb Z_K\ \in\ \bigl\langle\operatorname{im}\bigl(\Shv(T_0;D(\mathbb Z))
\xrightarrow{\ \pi_*\ }\Shv(K;D(\mathbb Z))\bigr)\bigr\rangle,
\end{equation*}
i.e.\ that $\pi_*\mathbb Z_{T_0}$ is a \emph{descendable} commutative algebra
object (\cref{prop:18:descendable}) in $\Shv(K;D(\mathbb Z))$. One is free to
choose the profinite cover, so it suffices to find one profinite cover for which
descendability holds. By pullback (and by K\"unneth) one reduces to
$K=[0,1]^N$, and it is enough to treat $K=[0,1]$.
\end{remark}

\begin{remark}[The Cantor cover of the interval and finite-dimensionality]
\label{rmk:20:cantor-cover}
For $K=[0,1]$ write the standard Cantor cover as an inverse limit of compact
Hausdorff spaces rather than of finite sets: at each stage take the disjoint
union of the two dyadic halves of $[0,1]$ (then of the four quarters, etc.),
each mapping onto $[0,1]$; the inverse limit is the Cantor set, surjecting onto
$[0,1]$. The pushforward from the Cantor set is the sequential colimit of the
pushforwards from these finite stages. There is a general fact that
descendability of a sequential colimit follows from descendability with a
\emph{uniform} exponent of nilpotence at each finite stage; at each stage one has
a Mayer--Vietoris cover of the closed pieces with only double (no triple)
intersections, giving exponent $\approx 2$, and the colimit $\approx 3$ or $4$.
For $[0,1]^N$ the exponent grows with $N$ but stays finite -- this is exactly
where finite-dimensionality is needed; for the Hilbert cube $[0,1]^{\mathbb N}$
the estimate degenerates and one can expect nothing.
\end{remark}

\begin{remark}[Two further remarks]
\label{rmk:20:two-remarks}
\leavevmode
\begin{itemize}
\item A posteriori (from the proof, not a priori), the condition in
\cref{thm:20:compact-hausdorff}(2) is equivalent to asking that the connectivity
estimate hold after a \emph{proper} cover of $\mathrm{Spec}(R)$, rather than only
$!$-locally.
\item If $Z\subseteq K$ is closed with open complement $U$, then the associated
analytic stacks are each other's complements in the naive way: the functor of
points of $U$ sends $R$ to those maps $\mathrm{Spec}(R)\to K$ whose pullback to
$Z$ is the empty analytic stack, and dually. (Check on profinite sets.)
\end{itemize}
The passage from $F((\pi_0)_*\mathbb Z)$ as an $\mathbb E_\infty$-algebra to an
\emph{animated} algebra requires an extra argument (derived symmetric powers;
work in progress with Maxime Ramzi, on a categorical characterisation of animated
commutative rings -- what structure on the module category of an $\mathbb
E_\infty$-ring promotes it to an animated ring, something like derived symmetric
power functors); Clausen deferred it. In characteristic $0$ it is automatic.
\end{remark}

\subsection{Norms on analytic rings}
\label{subsec:20:norms}

We now define norms. Scholze introduced the notion last time
(\cref{def:19:norm}); Clausen's axioms differ cosmetically (multiplicativity) and
substantively (Axiom~4). Fix an analytic ring $R$. Over any analytic ring one can
build the \emph{algebraic} projective line $\mathbb P^1_R$ (base change of
$\mathbb P^1_{\mathbb Z}$ with its trivial structure), and the closed interval
$[0,\infty]$ is a condensed set, hence an analytic stack.

\begin{definition}[Norm]
\label{def:20:norm}
A \emph{norm} on $R$ is a map of analytic stacks
\begin{equation*}
N\colon\ \mathbb P^1_R\ \longrightarrow\ [0,\infty]
\end{equation*}
subject to Axioms 1--4 below.
\end{definition}

Before listing the axioms, note what a norm does to a ring element. For $f\in
\ur R(\ast)$ one gets a section $\mathrm{Spec}(R)\to\mathbb A^1_R\subseteq
\mathbb P^1_R$, and composing with $N$ a map
\begin{equation*}
N(f)\colon\ \mathrm{Spec}(R)\ \longrightarrow\ [0,\infty].
\end{equation*}
So the norm of an element is \emph{not} a real number but a map
$\mathrm{Spec}(R)\to[0,\infty]$: think of it as a family of residue fields, with
a real number for each, possibly varying. The datum is geometric: given a map
$R\to R'$ of analytic rings, a norm on $R$ pulls back to a norm on $R'$ (by
composition), and persists.

\begin{itemize}
\item[\textbf{Axiom 1.}] $N(0)\colon\mathrm{Spec}(R)\to[0,\infty]$ factors
through $\{0\}$ (the analytic stack of the singleton point $\{0\}\subseteq
[0,\infty]$); i.e.\ the norm of $0$ is the constant function $0$.
\end{itemize}

\begin{remark}[The zero ring]
\label{rmk:20:zero-ring}
This geometric phrasing sidesteps the usual annoyance of whether $N(1)$ should be
$1$ or $0$ over the zero ring: for the zero ring $\mathrm{Spec}(R)=\emptyset$, so
$N(1)$ is simultaneously $1$ and $0$ (the empty map factors through both), and no
inconsistency arises.
\end{remark}

\begin{itemize}
\item[\textbf{Axiom 2 (compatibility with inversion).}] Writing $T$ for the
coordinate on $\mathbb P^1$, the square
\begin{equation*}
\begin{tikzcd}[column sep=large, row sep=large]
\mathbb P^1_R \arrow[r, "N"] \arrow[d, "T\mapsto T^{-1}"'] & {[0,\infty]}
\arrow[d, "\lambda\mapsto\lambda^{-1}"] \\
\mathbb P^1_R \arrow[r, "N"'] & {[0,\infty]}
\end{tikzcd}
\end{equation*}
commutes, both vertical maps exchanging $0$ and $\infty$.
\end{itemize}

Axioms~1 and~2 together imply that $N(\infty)$ factors through $\{\infty\}$, that
$N(1)$ factors through $\{1\}$, and (by the same reasoning) $N(-1)$ through
$\{1\}$.

\begin{definition}[Analytic affine line]
\label{def:20:analytic-affine-line}
Given a norm $N$ on $R$, the associated \emph{analytic affine line} is the open
locus where the norm is a finite nonzero real number,
\begin{equation*}
\Aan_R\ :=\ N^{-1}\bigl((0,\infty)\bigr)\ \subseteq\ \mathbb P^1_R .
\end{equation*}
\end{definition}

\begin{itemize}
\item[\textbf{Axiom 3 (multiplicativity).}] The square
\begin{equation*}
\begin{tikzcd}[column sep=huge, row sep=large]
\Aan_R\times\Aan_R \arrow[r, "N\times N"]
\arrow[d, "{(T,S)\mapsto TS}"'] & {(0,\infty)\times(0,\infty)}
\arrow[d, "{(\lambda,\lambda')\mapsto\lambda\lambda'}"] \\
\Aan_R \arrow[r, "N"'] & {[0,\infty]}
\end{tikzcd}
\end{equation*}
commutes; i.e.\ $N(TS)=N(T)N(S)$ on the analytic affine line.
\end{itemize}

The multiplication $(T,S)\mapsto TS$ needs justification: it is a priori defined
only on the algebraic affine line, and we must know it restricts to $\Aan_R$.

\begin{lemma}[Multiplication is defined on $\Aan_R$]
\label{lem:20:mult-well-defined}
$\Aan_R\subseteq\mathbb P^1_R$ is contained in the algebraic affine line
$\mathbb A^1_R$ (recall $D(\mathbb A^1_R)=\Mod_{R[T]}(D(R))$), and multiplication
carries $\Aan_R\times\Aan_R$ into it.
\end{lemma}

\begin{proof}[Sketch]
Since $N(\infty)=\infty$ (Axioms 1--2), the $\infty$-section of $\mathbb P^1$ is
an algebra over $N^*(\mathbb Z_{\{\infty\}})$, so $\Aan_R$ misses the
$\infty$-section. It remains to see that a stack over $\mathbb A^1$ missing the
$0$-section factors through $\mathbb G_m$; this is a special case of the
following general claim.

\emph{Claim.} If $X$ is an analytic stack over $\mathbb A^1_R$ whose pullback to
the $0$-section is the empty stack, then $X\to\mathbb A^1_R$ factors through
$(\mathbb G_m)_R$. Indeed, one may assume $X=\mathrm{Spec}(A)$ affine and think of
the map via the pullback functor $f^*\colon D(\mathbb A^1_R)\to D(A)$. Missing
the $0$-section means $f^*$ kills the structure sheaf of the origin
$\mathcal O_{\mathbb A^1}/T$; hence it kills everything built from it under
colimits, so it does not distinguish an object from the result of inverting $T$,
i.e.\ $f^*$ factors through $D\bigl((\mathbb G_m)_R\bigr)$. \qedhere
\end{proof}

For the fourth axiom we need the ``basic disk''. Recall from
\cref{lec:5,lec:14} the free complete module on a null sequence,
\begin{equation*}
P\ =\ R[\mathbb N\cup\{\infty\}]/R[\infty],
\end{equation*}
which carries a ring structure (measures on $\mathbb N$, multiplication from
addition on $\mathbb N$; \cref{rmk:7:P-ring}) and a ring map $R[T]\to P_R$
sending $T$ to the Dirac measure at $1$. Set
\begin{equation*}
D\ :=\ \mathrm{Spec}(P),\qquad D_R\ :=\ \mathrm{Spec}(P_R).
\end{equation*}
The map $R[T]\to P_R$ gives $D_R\to\mathbb A^1_R$; it is \emph{not} a
monomorphism (so $D_R$ is not literally a subset of $\mathbb A^1_R$), but it is
proper (hence $!$-able) -- a genuine disk that is not a subspace.

\begin{itemize}
\item[\textbf{Axiom 4.}] With $D_R\to\mathbb P^1_R$ as above:
\begin{enumerate}
\item[(4a)] the composite $D_R\to\mathbb P^1_R\xrightarrow{N}[0,\infty]$ factors
through $[0,1]\subseteq[0,\infty]$; and
\item[(4b)] the map $D_R\to\mathbb P^1_R$ is an isomorphism over the open unit
disk $N^{-1}\bigl([0,1)\bigr)$; i.e.\ in the square
\begin{equation*}
\begin{tikzcd}[column sep=large]
D_R \arrow[r] \arrow[rd] & \mathbb P^1_R \arrow[d, "N"] \\
& {[0,\infty]}
\end{tikzcd}
\qquad
\begin{tikzcd}[column sep=large, row sep=large]
N^{-1}\bigl([0,1)\bigr) \arrow[r, hook] \arrow[d] & \mathbb P^1_R
\arrow[d, "N"] \\
{[0,1)} \arrow[r, hook] & {[0,\infty]}
\end{tikzcd}
\end{equation*}
the pullback of $D_R$ over $N^{-1}([0,1))$ is $N^{-1}([0,1))$ itself.
\end{enumerate}
\end{itemize}

So $D_R$ is ``sandwiched between the closed and the open unit disk'': its norm is
$\le 1$, and it agrees with $\mathbb P^1_R$ (is a genuine subspace) over the open
unit disk, while blowing up (failing to be a monomorphism) along the boundary
$N=1$.

\begin{remark}[Motivation for Axiom 4: $P_R$ as a disk between open and closed]
\label{rmk:20:disk-motivation}
The ring map $R[T]\to P_R$ induces an isomorphism modulo $T^n$ for every $n$, so
$D_R\to\mathbb A^1_R$ is an isomorphism on the formal neighbourhood of the origin
and blows up as one moves away from it. In most examples $P_R$ lives in the heart
of $D(R)$ and one has injections
\begin{equation*}
R[T]\ \hookrightarrow\ P_R\ \hookrightarrow\ R[[T]],
\end{equation*}
so that $P_R$ is a space of sequences
\begin{equation*}
P_R\ =\ \bigl\{(r_0,r_1,\dots)\ :\ \text{termwise multiplication by a null
sequence yields a ``summable'' sequence}\bigr\},
\end{equation*}
where the notion of summability depends only on the analytic ring structure. (For
intuition one may take $R=\mathbb R$ and ``summable'' $=$ absolutely summable --
not literally a special case, but close enough.) This is forced by the universal
property: maps $P_R\to M$ in $D(R)$ correspond to null sequences in $M$, paired
by summation. For $R=\mathbb C$ with its gaseous or liquid structure
(\cref{def:12:alpha-liquid-R}), one gets
\begin{equation*}
\mathcal O^{\mathrm{hol}}\bigl(\{|z|\le1\}\bigr)\ \longrightarrow\ P_{\mathbb C}\
\longrightarrow\ \mathcal O^{\mathrm{hol}}\bigl(\{|z|<1\}\bigr),
\end{equation*}
the summability condition being weaker than overconvergence on the closed unit
disk (a power series with a value at $1$ has summable coefficients) but stronger
than convergence on the open disk. So $P$ sits strictly between the open and
closed unit disk of radius $1$; Axiom~4 axiomatises exactly this placement.
\end{remark}

\begin{remark}[Status of Axiom 4]
\label{rmk:20:axiom4-status}
Axiom~4 was not among Scholze's conditions (\cref{def:19:norm}); Clausen adds it,
and it is not known whether it follows from Axioms 1--3. It \emph{does} follow
when the base solidifies to $0$ -- e.g.\ if $\mathrm{Spec}(R)$ base-changed to
$\mathrm{Spec}(\mathbb Z_\solid)$ is the empty analytic stack, as happens for any
$R$ whose underlying ring is an algebra over $\ur{\mathbb R}$. Non-archimedean
Huber-ring norms are not solid, so there could a priori be pathological norms
violating Axiom~4 there; but the usual norms all satisfy it, and Clausen takes it
as an axiom. There is no additivity or triangle-inequality condition anywhere --
$N$ may be archimedean or non-archimedean at will.
\end{remark}

\begin{remark}[What a norm supplies: overconvergent disks]
\label{rmk:20:overconvergent-disks}
For $r\in\mathbb R_{\ge0}$, the pullback $N^*(\mathbb Z_{[0,r]})$ is an idempotent
algebra in $D(\mathbb P^1_R)$, to be interpreted as \emph{overconvergent}
functions on the closed disk of radius $r$. Overconvergence is forced: the
constant sheaf on $[0,r]$ is the filtered colimit
\begin{equation*}
\mathbb Z_{[0,r]}\ =\ \varinjlim_{r'>r}\mathbb Z_{[0,r']},
\end{equation*}
and $N^*$ (a pullback) commutes with colimits, so any version of ``functions on
the disk of radius $r$'' receiving a map from these is forced to be the
overconvergent one. The data of these idempotent algebras, together with the
inversion Axiom~2, determines $N$: the closed subsets of $[0,\infty]$ are easy to
classify, one need only give the loci $\{N\le r\}$ and (by inversion)
$\{N\ge r\}$, subject to simple compatibilities.
\end{remark}

\begin{remark}[Why norms, not just a base ring]
\label{rmk:20:why-norms}
Over a bare analytic ring the only geometry one knows how to do is to import
algebraic geometry (\cref{prop:18:zariski}). A norm supplies the extra structure
-- disks of prescribed radius inside the affine line, which turn out to be
$!$-able subsets -- needed to do something resembling traditional analytic
geometry (open and closed disks of a given radius). Without a norm there is no
notion of the radius of a disk.
\end{remark}

\subsection{Classifying norms}
\label{subsec:20:classifying}

We now classify norms on a given analytic ring $R$. The first step lets us assume
a topologically-small unit exists.

\begin{lemma}[A small element exists locally]
\label{lem:20:small-element}
Given a norm $N\colon\mathbb P^1_R\to[0,\infty]$, the locus $N^{-1}\bigl((0,1)
\bigr)\to\mathrm{Spec}(R)$ is a cover in the $!$-topology. Equivalently, working
locally on $\mathrm{Spec}(R)$ we may assume there is an element $q\in\ur R(\ast)$
with $N(q)\in(0,1)$.
\end{lemma}

\begin{proof}[Sketch]
It suffices, and is stronger, to show $N^{-1}(\{\tfrac12\})\to\mathrm{Spec}(R)$
is a cover. For simplicity assume $N^*(\mathbb Z_{\{1/2\}})$ is connective, so
that $N^{-1}(\{\tfrac12\})$ is affine, and proper over $\mathbb A^1_R$, hence over
$\mathrm{Spec}(R)$ (in general one refines by a proper descendable cover by an
affine and combines the generation statements; we omit this). By the
$!$-cover criterion for proper maps (\cref{prop:18:check}(2a)) we must see
$\ur R\in\langle N^*(\mathbb Z_{\{1/2\}})\rangle$.

\emph{Claim: $\ur R$ is a direct retract of $N^*(\mathbb Z_{\{1/2\}})$.} In
$D(\mathbb P^1_R)$ the constant sheaf on $[0,\infty]$ is the pullback (equivalently
pushout) of the constant sheaves on $[0,\tfrac12]$ and $[\tfrac12,\infty]$ glued
along the point $\{\tfrac12\}$:
\begin{equation*}
\begin{tikzcd}[column sep=large, row sep=large]
\mathcal O_{\mathbb P^1_R} \arrow[r] \arrow[d]
& N^*(\mathbb Z_{[0,1/2]}) \arrow[d] \\
N^*(\mathbb Z_{[1/2,\infty]}) \arrow[r] & N^*(\mathbb Z_{\{1/2\}}).
\end{tikzcd}
\end{equation*}
Pushing forward to $D(R)$, the cohomology of $\mathcal O_{\mathbb P^1_R}$ is
$\ur R$ in degree $0$, and the pushout presents $\ur R$ as a homotopy pushout of
the three pushed-forward corners. To retract $\ur R$ off $N^*(\mathbb Z_{\{1/2\}})$
one builds a map $\ur R\to N^*(\mathbb Z_{\{1/2\}})$ using the two evaluations at
$0$ and $\infty$ (each landing in $\ur R$), which agree on restriction, so give a
map splitting the unit -- concretely, picking out the constant term of a Laurent
series. Hence $\ur R$ is a retract of $N^*(\mathbb Z_{\{1/2\}})$, which is thus
descendable of index $0$: the descent is direct, by a retract.
\end{proof}

The next result says a topologically nilpotent unit visible to the norm must come
from the gaseous base ring of \cref{lec:14},
\begin{equation*}
\mathbb Z[\widehat q^{\pm1}]_\gasx,
\end{equation*}
the initial analytic ring with a topologically nilpotent unit $q$ for which
$1-q\cdot\mathrm{shift}$ is invertible on $P\otimes_{\mathbb Z}(-)$
(\cref{thm:14:gaseous-base,rmk:19:gaseous-base}).

\begin{proposition}[Small elements are gaseous]
\label{prop:20:gaseous-factorization}
Given a norm on $R$ and $q\in\ur R(\ast)$ with $N(q)\in(0,1)$, the ring map
$\mathbb Z[T]\to\ur R$, $T\mapsto q$, factors uniquely through
$\mathbb Z[\widehat q^{\pm1}]_\gasx$.
\end{proposition}

\begin{proof}[Sketch]
\emph{Uniqueness.} This has nothing to do with norms: it is the claim that
\begin{equation*}
\mathrm{Spec}\bigl(\mathbb Z[\widehat q^{\pm1}]_\gasx\bigr)\ \longrightarrow\
\mathbb A^1_{\mathbb Z}
\end{equation*}
is a monomorphism (so the space of factorisations is contractible if nonempty).
This is not obvious from the construction of the gaseous ring, which inverted $q$
and imposed gaseousness on $\mathbb Z[\widehat q]=P$, an object that is
\emph{not} idempotent. Instead present
$\mathbb Z[\widehat q^{\pm1}]_\gasx=\Mod_P\bigl(\mathbb Z[q^{\pm1}]_\gasx\bigr)$,
i.e.\ view it as $P$-modules in the polynomial ring with a \emph{gaseous}
variable $q$. Then $P$ \emph{is} idempotent in $D(\mathbb Z[q^{\pm1}]_\gasx)$:
\begin{equation*}
P\otimes^{}_{}P\ \xrightarrow{\ \sim\ }\ P
\quad\text{in }D\bigl(\mathbb Z[q^{\pm1}]_\gasx\bigr),
\end{equation*}
which after base change reduces to the computation
\begin{equation*}
P_{\mathbb Z[q^{\pm1}]_\gasx}\big/(T-q)\ =\ \mathbb Z[q^{\pm1}]_\gasx
\end{equation*}
(here $T$ is the variable coming from $P$ and $q$ the unit); modding out by
$T-q$ returns the underlying ring, checked via a short exact sequence in the
free-module description of $P$.

\emph{Existence.} From $N(q)\in(0,1)$: being in $(0,1)$ means being away from
$N^{-1}([1,\infty])$, hence away from the ``disk $D$ sitting at $\infty$''
(the image of $D$ under inversion), which by \cref{lem:20:mult-well-defined}'s
argument forces $q$ to be gaseous; and being in $[0,1)$ forces $q$ topologically
nilpotent (it comes from $D$, by Axiom~4). Together these say $q$ factors through
$\mathbb Z[\widehat q^{\pm1}]_\gasx$. \qedhere
\end{proof}

\begin{theorem}[Classification of norms]
\label{thm:20:classification}
Let $R$ be an analytic ring. The map recording $q$ and its norm,
\begin{equation}
\label{eq:20:classification}
\left\{\begin{gathered}N\colon\mathbb P^1_R\to[0,\infty]\ \text{a norm on }R,\\
q\in\ur R(\ast)\ \text{with}\ N(q)\in(0,1)\end{gathered}\right\}
\ \xrightarrow{\ \ \sim\ \ }\
\left\{\begin{gathered}\mathrm{Spec}(R)\xrightarrow{(q,\,N(q))}\\
\mathrm{Spec}\bigl(\mathbb Z[\widehat q^{\pm1}]_\gasx\bigr)\times(0,1)\end{gathered}\right\}
\end{equation}
is an equivalence.
\end{theorem}

So, up to a cover of $\mathrm{Spec}(R)$, giving a norm on $R$ is the same as
giving, for a chosen element $q$: the value $N(q)\colon\mathrm{Spec}(R)\to(0,1)$,
and a witness that $q$ comes from the gaseous theory (the map to
$\mathrm{Spec}(\mathbb Z[\widehat q^{\pm1}]_\gasx)$).

\begin{remark}[Both sides are ``sets'' relative to $\mathbb A^1$]
\label{rmk:20:relative-sets}
Both sides of \eqref{eq:20:classification} map to $\mathbb A^1$ (recording $q$),
with fibres that are sets; they are ``as much sets as each other'' but neither is
individually a set, since $q$ is a derived datum (chosen in $\pi_0\bigl(\ur R
(\ast)\bigr)$, with the higher structure coming along automatically or
requiring construction).
\end{remark}

\begin{proof}[Sketch]
After a cover of $\mathrm{Spec}(R)$ one may assume $q$ admits compatible $n$-th
roots for all $n$: the map $\mathbb Z[T]\to\mathbb Z[T^{1/n}:\text{all }n]$ is
fpqc and countably presented, hence descendable (\cref{ex:18:descendable}), and
both sides of \eqref{eq:20:classification} are sheaves, so we may work locally.
Rescaling the norm we may also assume $N(q)\equiv\tfrac12$; then one checks
$N(q^{1/n})=(\tfrac12)^{1/n}$ (one inclusion by taking $n$-th powers, the other by
passing to complementary opens), so in particular roots of unity have norm $1$.

For $\alpha\in\mathbb G_m$ write $D_\alpha$ for the pullback of $D\to
\mathbb A^1_R$ along multiplication by $\alpha$, $\mathbb A^1_R\xrightarrow{\cdot
\alpha}\mathbb A^1_R$; this is a disk ``of radius $\alpha$''. If $\alpha=q\beta$
with $q$ topologically nilpotent, one gets a map $D_\alpha\xrightarrow{\cdot q}
D_\beta$ (using that $P$ is a Hopf algebra, encoding multiplication) -- an
inclusion of disks. Applying this to $\alpha=q^{m/n}$, $m/n\in\mathbb Q$, produces
disks of radius $N(q)^{m/n}$ with inclusions between them; taking overconvergent
versions, the blow-up-along-the-boundary defects wash out in the colimit and one
recovers exactly the datum a norm must supply. This shows the map is a
monomorphism; conversely, running the construction (translate the fake disk by
powers of $q$, form the inclusions, make them overconvergent, check the axioms)
produces a norm from the right-hand data, giving surjectivity.
\end{proof}

\begin{remark}[Relation to Berkovich norms; and what norms are for]
\label{rmk:20:berkovich-relation}
The assignment $R\mapsto\{\text{norms on }R\}$ is itself an analytic stack (it
fits the definition, and \cref{lem:20:small-element} exhibits a cover). Only
Tate (analytic) rings admit nontrivial norms: a discrete ring has no nontrivial
theory of norms. For a Tate--Huber ring, a norm in this sense amounts to a
choice of $|q|\in(0,1)$ for a pseudo-uniformizer $q$, yielding a
non-archimedean norm on each residue field, chosen continuously; the space of
norms is a torsor under rescaling ($N\mapsto N^\alpha$) of a fixed one, matching
the rescaling on the Berkovich spectrum. For a global object this is far from a
single Berkovich seminorm: it is a whole compatible family. Modding out by these
exponential rescalings (so that one can no longer speak of a disk of a
\emph{fixed} radius, but much else survives) gives a larger stack that also
accommodates $\mathbb Z_\solid$ and objects over it; this is to be discussed
later.
\end{remark}

\begin{remark}[Reconstruction from normed points]
\label{rmk:20:reconstruction}
By analogy with the theory of diamonds -- where a discrete (or general) ring is
substantially recovered from how it maps to perfectoid Tate--Huber rings, and
Scholze remembers only such maps -- one can restrict attention to normed
analytic rings and recover an object essentially from its functor of points on
normed analytic rings alone.
\end{remark}
